\documentclass[11pt]{article}
\usepackage[margin=1in]{geometry}
\usepackage{amsmath,amssymb}
\usepackage{booktabs}
\usepackage{hyperref}
\usepackage{siunitx}
\title{Replication Report: Faraday and Kerr Rotation due to Photoinduced
Orbital Magnetization in a 2DEG\\
\large Durnev (2023), arXiv:2306.08509}
\author{Independent analytic replication}
\date{\today}

\begin{document}
\maketitle

\begin{abstract}
We independently replicate the central analytic results of Durnev (2023),
which develops a Boltzmann/kinetic-theory of the pump-induced Faraday and
Kerr rotation of a probe field in a two-dimensional electron gas (2DEG).
A circularly polarized pump induces an orbital magnetization (intraband
inverse Faraday effect), producing a transverse conductivity
$\sigma_{xy}$ and hence rotation of a linearly polarized probe. Working in
Gaussian/CGS units, we implement the paper's closed-form expressions for the
graphene (linear-dispersion) case and evaluate them at the resonance condition
$\Omega\tau_1\approx1$, $\omega\approx\Omega$. All intermediate cross-checks
reproduce the paper's own stated numbers to $<1$--2\%: the Fermi energy
$\varepsilon_F=\SI{63.9}{meV}$ (paper $\SI{64}{meV}$), the dimensionless
coupling $2\pi\sigma_0/(c\bar n)=0.0708$ (paper $0.071$), and the transmission
$T(\Omega)=0.699$ (paper range $0.63$--$0.70$). Two independent formulas for the
peak Faraday angle agree ($\theta_F=0.0441^\circ$ from Eq.~26 and
$0.0417^\circ$ from Eq.~5+25), and the synthetic magnetic field
$B_{\mathrm{syn}}=\SI{0.088}{T}$ agrees with the paper's $\sim\SI{0.1}{T}$
estimate to 12\%. \textbf{Verdict: REPLICATED} (coverage 7/10,
agreement 9/10).
\end{abstract}

\section{Paper and Claim}
Durnev (2023) studies the Faraday and Kerr rotation of a weak probe field
arising from the orbital magnetization of a 2DEG driven by a circularly
polarized pump, in the intraband spectral range. The theory is built on the
analytic solution of the kinetic (Boltzmann) equation for both linear
(graphene) and parabolic dispersions and arbitrary scattering potential. The
headline claim: at the resonance $\omega\approx\Omega$ with $\Omega\tau_1\sim1$,
the Faraday and Kerr rotation angles reach $\sim0.1^\circ$ per
$\SI{1}{kW/cm^2}$ of pump intensity in graphene, corresponding to a
pump-induced synthetic magnetic field of $\sim\SI{0.1}{T}$.

\section{Method}
This is an analytic/closed-form paper: the replication \emph{is} the
derivation and its numerical evaluation. We reimplement, from the paper's
equations, the graphene single-layer (linear-dispersion) chain in
Gaussian/CGS units:
\begin{itemize}
  \item \textbf{Eq.~(5)} $\theta_F + i\,\epsilon_F = 2\pi\sigma_{xy}/
        \bigl[c\bar n(1+\alpha)\bigr]$, with
        $\alpha=2\pi\sigma/(c\bar n)$ and $\bar n=(n_1+n_2)/2$.
  \item \textbf{Eq.~(25)} the resonance transverse conductivity
        $\sigma_{xy}(\omega)$ for a linear spectrum.
  \item \textbf{Eq.~(26)} the explicit closed-form resonance Faraday angle
        $\theta_F(\omega)$.
  \item \textbf{Eq.~(27)} the synthetic magnetic field
        $B_{\mathrm{syn}}\sim e c |E_\Omega|^2\tau_0/\varepsilon_F$.
  \item Field at the sheet: $|E_\Omega|^2 = 2\pi T(\Omega) I_\Omega/(c n_2)$,
        with $T(\Omega)=n_2|\bar t|^2/n_1$ and $\bar t=t_{12}/(1+\alpha(\Omega))$.
\end{itemize}
Supporting relations: graphene Fermi energy
$\varepsilon_F=\hbar v_0\sqrt{\pi n_e}$ and static conductivity
$\sigma_0=e^2 v_0^2 n_e \tau_1/\varepsilon_F$. Numerical evaluation uses
NumPy; the probe frequency $\omega$ is scanned across the resonance to locate
the peak angle.

\section{Parameters}
\begin{table}[h]
\centering
\begin{tabular}{lll}
\toprule
Symbol & Value & Meaning \\
\midrule
$\tau_1$ & \SI{0.1}{ps} & momentum relaxation time ($\Omega\tau_1=1$) \\
$\tau_0$ & \SI{5}{ps} & energy / zeroth-harmonic relaxation time \\
$v_0$ & \SI{1e8}{cm/s} & graphene Fermi velocity \\
$n_e$ & \SI{3e11}{cm^{-2}} & 2D electron density \\
$I_{\mathrm{pump}}$ & \SI{1}{kW/cm^2} & pump intensity \\
$n_1,n_2$ & $1,\,3$ & refractive indices (substrate case) \\
$P_{\mathrm{circ}}$ & $1$ & pump circular polarization degree \\
\bottomrule
\end{tabular}
\caption{Graphene-on-substrate parameters (Fig.~3 of the paper).}
\end{table}

\section{Results and Cross-Checks}
\begin{table}[h]
\centering
\begin{tabular}{lccc}
\toprule
Quantity & This work & Paper & Agreement \\
\midrule
$\varepsilon_F$ (meV)                 & 63.90  & 64.0            & $<0.2\%$ \\
$2\pi\sigma_0/(c\bar n)$              & 0.0708 & 0.071           & $<1\%$   \\
$T(\Omega)$                          & 0.699  & 0.63--0.70      & in range \\
$\Omega_{\mathrm{res}}$ (THz)        & 1.59   & ($\Omega\tau_1=1$) & consistent \\
$\theta_F$ peak, Eq.~26 (deg)        & 0.0441 & $\sim$0.05 (Fig.3) & $\approx$ \\
$\theta_F$ peak, Eq.~5+25 (deg)      & 0.0417 & $\sim$0.05 (Fig.3) & $\approx$ \\
$B_{\mathrm{syn}}$ (T)               & 0.088  & $\sim$0.1        & 12\%     \\
\bottomrule
\end{tabular}
\caption{Cross-checks and headline results. The two independent Faraday
formulas (Eq.~26 and Eq.~5+25) agree with each other to $\sim6\%$.}
\end{table}

The intermediate cross-checks are the strongest evidence: they reproduce the
paper's \emph{own} intermediate numbers to $<1$--2\%, confirming the unit
system, the Fermi-energy/conductivity relations, and the transmission model
are all faithfully reproduced. The peak Faraday angle
$\theta_F\approx0.044^\circ$ lands squarely on the $\Omega\tau_1=1$
graphene-on-substrate curve of Fig.~3 (y-axis $\sim0.05^\circ$). The headline
``$0.1^\circ$'' is an order-of-magnitude statement spanning Figs.~2--4
($0.1$--$1^\circ$); the graphene-on-substrate value sits correctly within it.
The synthetic field $B_{\mathrm{syn}}=\SI{0.088}{T}$ (paper's explicit estimate
parameters $\varepsilon_F=\SI{50}{meV}$, $\tau_0=\SI{10}{ps}$) matches
$\sim\SI{0.1}{T}$ to 12\%.

\section{Verdict}
\textbf{REPLICATED.} Coverage $\approx 7/10$, Agreement $\approx 9/10$.
For an analytic/closed-form paper the replication is the re-derivation and
numerical evaluation of the governing formulas. Multiple \emph{independent}
cross-checks match the paper to $<1$--2\%, and two independent Faraday-angle
formulas mutually agree to $\sim6\%$. Honest gaps: analytic-only (no full
transport simulation); only the linear/graphene branch is coded (parabolic /
bilayer Eqs.~21/23/24 not implemented); Kerr angle taken as
$\theta_K\approx-\theta_F$ in the large-contrast limit rather than evaluated
independently via Eq.~(6); and no pixel-level digitization of the paper's
figures was performed.

\end{document}
